% ECONOMICS_PROBLEM: {"id": "C72-4", "title": "Computational complexity of optimin in finite games", "jel_code": "C72", "source_id": "OP-0589"}
% !TeX program = xelatex
\documentclass[11pt,a4paper]{article}
\usepackage[margin=23mm]{geometry}
\usepackage{fontspec}
\setmainfont{Latin Modern Roman}
\usepackage{amsmath,amssymb,mathtools}
\usepackage{xcolor,enumitem,booktabs,longtable,array,xurl,hyperref}
\definecolor{muted}{HTML}{53636C}
\hypersetup{colorlinks=true,urlcolor=blue,linkcolor=blue}
\setlength{\parindent}{0pt}
\setlength{\parskip}{5pt}
\newcommand{\R}{\mathbb R}
\newcommand{\E}{\mathbb E}
\newcommand{\N}{\mathbb N}
\newcommand{\ind}{\mathbf 1}
\newcommand{\Var}{\operatorname{Var}}
\newcommand{\Cov}{\operatorname{Cov}}
\newcommand{\supp}{\operatorname{supp}}
\DeclareMathOperator*{\argmin}{arg\,min}
\DeclareMathOperator*{\argmax}{arg\,max}
\newcommand{\entrymeta}[3]{{\sffamily\small\color{muted}\textbf{\texttt{#1}}\quad|\quad #2\par #3\par}}
\newcommand{\Problemref}[1]{\texttt{#1}}
\newcommand{\JELref}[2]{\texttt{#2} (JEL \texttt{#1})}
\begin{document}
\section*{Computational complexity of optimin in finite games}
\textbf{Problem ID:} \texttt{C72-4}\quad \textbf{JEL:} \texttt{C72}\par
% BEGIN ECONOMICS STATEMENT
\label{problem:OP-0589}
\noindent\textbf{Addendum ID:} OPT-COMP-1. \textbf{Formulation:} editorial family of precise computational problems motivated by Ismail (2025), not a conjecture attributed verbatim to that article.

\subsection*{Definitions and assumptions}

Let $G=(N,(A_i)_{i\in N},(v_i)_{i\in N})$ be a finite normal-form game, with $N=\{1,\ldots,n\}$, $n\ge2$, $m_i=|A_i|\ge2$, $A=\prod_i A_i$, and rational payoffs $v_i:A\to[-1,1]\cap\mathbb Q$. Set $\Delta_i=\Delta(A_i)$ and $\Delta=\prod_i\Delta_i$. Write $u_i(p)$ for the multilinear expected-payoff extension of $v_i$ at $p\in\Delta$. The default encoding is the explicit table of the $n|A|$ payoff entries in binary, of total bit length $L$. All complexity assertions refer to $L$ (and to the bit length of any additional rational inputs), not merely to the number of pure actions.

For each player $j$ and mixed profile $p$, define the correspondence of permitted \emph{strictly profitable unilateral deviations from the baseline $p$}, including non-deviation, by
\[
B_j(p)=\{p_j\}\ \cup\
\bigl\{q_j\in\Delta_j:
u_j(q_j,p_{-j})>u_j(p)\bigr\},
\qquad B_{-i}(p)=\prod_{j\in N\setminus\{i\}}B_j(p).
\]
The \emph{optimin performance} of player $i$ and the mixed-strategy optimin correspondence are, respectively,
\[
\pi_i^G(p)=\inf_{q_{-i}\in B_{-i}(p)}u_i(p_i,q_{-i}),
\]
\[
\begin{aligned}
\mathcal O^{\mathrm{mix}}(G)=\bigl\{p\in\Delta:\;&\nexists q\in\Delta\text{ such that }
\pi_i^G(q)\ge\pi_i^G(p)\ \forall i\in N,\\
&\hspace{21mm}\text{and }
\pi_j^G(q)>\pi_j^G(p)\text{ for some }j\in N\bigr\}.
\end{aligned}
\]
Each $q_j\in B_j(p)$ is assessed against the \emph{original} profile $p_{-j}$, even when several opponents deviate in the product $B_{-i}(p)$. Thus the definition permits combinations of independently profitable unilateral deviations; it does not require the combination to be jointly profitable. In particular, a Nash equilibrium $p$ has $B_j(p)=\{p_j\}$ for all $j$. The strict inequality in $B_j(p)$ must not silently be replaced by a weak inequality.

The \emph{pure-restricted} variant is defined separately: replace $\Delta_j$ by $A_j$ in every deviation set and in the set of candidate profiles, denote the resulting performance by $\pi_i^{G,\mathrm{pure}}(a)$, and let $\mathcal O^{\mathrm{pure}}(G)\subseteq A$ consist of the profiles Pareto undominated under $\pi^{G,\mathrm{pure}}$. Membership in $\mathcal O^{\mathrm{pure}}(G)$ does not imply membership in $\mathcal O^{\mathrm{mix}}(G)$.

For exact mixed-strategy search, specify a finite algebraic-number representation of each output probability, or an explicitly equivalent exact real-algebraic representation. For approximate search, output rational probabilities. The required error notion and its binary-encoded rational tolerance $\varepsilon>0$ are always stated explicitly.

\subsection*{Formal research question}

\textbf{Central problem.} Determine the exact worst-case computational complexity of constructing one element of $\mathcal O^{\mathrm{mix}}(G)$ for a rational finite game $G$. Classify the following interrelated, but distinct, search, decision, verification, and approximation problems. Any proposed completeness claim must specify its many-one, Turing, or search-reduction notion and its output encoding.

\begin{enumerate}[label=\textbf{(\Alph*)},leftmargin=2.7em]

\item \textbf{Two-player general-sum mixed optimin.}
For the input class $n=2$ with arbitrary rational bimatrix payoffs, determine the complexity of the total search task
\[
\mathsf{MIXED\text{-}OPTIMIN}_2:\quad
G\longmapsto p\in\mathcal O^{\mathrm{mix}}(G).
\]
Does an algorithm run in time polynomial in the bimatrix input length $L$ under a suitable exact output model? Otherwise, establish a sharp hardness or completeness classification, distinguishing arbitrary $m_1,m_2$ from bounded-action subfamilies. Investigate membership or hardness relative to appropriately defined total-search classes such as $\mathsf{PPAD}$ and exact real-algebraic search classes such as $\mathsf{FIXP}$; neither classification is assumed.

\item \textbf{Three-player and general $n$-player mixed optimin.}
For each fixed $n\ge3$, and separately when $n$ is part of the input, determine the search complexity of
\[
\mathsf{MIXED\text{-}OPTIMIN}_n:\quad
G\longmapsto p\in\mathcal O^{\mathrm{mix}}(G).
\]
Identify the least player count at which the complexity provably differs from the two-player problem, if such a separation exists. Distinguish general-sum games from $n$-player constant-sum games ($\sum_i v_i(a)=c$ for every $a\in A$), and distinguish a fixed number of actions per player from unrestricted action counts. In particular, determine whether the three-player problem admits a $\mathsf{FIXP}$-type characterization, a different real-algebraic search characterization, or a strictly simpler algorithmic description; no transfer from Nash-equilibrium complexity is presumed.

\item \textbf{Approximation by strategy distance versus performance efficiency.}
Let $d=\sum_i m_i$, and use the sup norm on the coordinates of $\Delta\subseteq\mathbb R^d$. Classify the total rational-output approximation problem
\[
\mathsf{NEAR\text{-}OPTIMIN}_n:\quad
(G,\varepsilon)\longmapsto\widehat p\in\Delta\cap\mathbb Q^d
\quad\text{with}\quad
\inf_{p^*\in\mathcal O^{\mathrm{mix}}(G)}
\|\widehat p-p^*\|_\infty\le\varepsilon.
\]
Separately, investigate the \emph{performance-based} requirement that there be no $q\in\Delta$ for which
\[
\pi_i^G(q)\ge\pi_i^G(\widehat p)\quad\forall i,
\qquad \max_i\bigl(\pi_i^G(q)-\pi_i^G(\widehat p)\bigr)>\varepsilon.
\]
Determine the conditions under which such a performance-based approximation admits a rational output and, when it does, its complexity. Compare polynomial dependence on $L$ and $\log(1/\varepsilon)$ with polynomial dependence on $L$ and $1/\varepsilon$; do not conflate closeness in strategies with closeness in performance.

\item \textbf{Performance thresholds and scalarized optimization.}
For $r\in([-1,1]\cap\mathbb Q)^n$, determine the decision complexity of
\[
\mathsf{OPTIMIN\text{-}GUARANTEE}_n(G,r):\quad
\exists p\in\Delta\ \forall i\in N:\ \pi_i^G(p)\ge r_i\,?
\]
For rational weights $\lambda_i>0$ satisfying $\sum_i\lambda_i=1$, define
\[
V_\lambda(G)=\max_{p\in\Delta}\sum_{i=1}^n\lambda_i\pi_i^G(p).
\]
Classify deciding $V_\lambda(G)\ge t$ for a rational threshold $t$, computing $V_\lambda(G)$ exactly or within $\varepsilon$, and constructing a maximizing profile. Every exact maximizer for strictly positive weights is an optimin point, but varying $\lambda$ need not generate every Pareto-undominated performance vector. Determine sharp complexity bounds, including those based on quantified polynomial inequalities over the reals, for these precisely formulated tasks.

\item \textbf{Checking and certifying a proposed mixed optimin.}
Given $G$ and a rational profile $p\in\Delta$, classify the decision problem
\[
\mathsf{IS\text{-}OPTIMIN}_n(G,p):\quad p\in\mathcal O^{\mathrm{mix}}(G)\,?
\]
Also classify the feasibility of polynomial-size, polynomial-time-verifiable certificates of optimin membership or nonmembership under declared exact and approximate encodings. Distinguish this global Pareto-nondominance test from evaluation of the individual numbers $\pi_i^G(p)$ and from checking ordinary Nash inequalities.

\item \textbf{Succinct input and the pure-strategy boundary.}
If the game is given as an \emph{explicit payoff table}, computing some $a\in\mathcal O^{\mathrm{pure}}(G)$ is already polynomial-time solvable, even with variable $n$; this is a settled baseline, not an open complexity classification. Instead, let the game be encoded by a Boolean or arithmetic circuit $C$ returning the rational payoff $v_i(a)$ given a player index $i$ and a pure action profile $a$. Determine the complexity of finding $a\in\mathcal O^{\mathrm{pure}}(G_C)$ and of deciding
\[
\exists a\in\mathcal O^{\mathrm{pure}}(G_C)
\quad\forall i:\quad \pi_i^{G_C,\mathrm{pure}}(a)\ge r_i,
\]
particularly when every player has two actions and the number of players grows. Record the circuit and payoff bit lengths, distinguish succinct from tabular input, and examine bounded-degree graphical or other compact representations separately.

\item \textbf{Structural restrictions and parameterized complexity.}
For every stated game representation, characterize tractable and hard subfamilies among two-person zero-sum, $n$-person constant-sum ($n\ge3$), common-interest, potential, symmetric, polymatrix, anonymous, and bounded-interaction graphical games. Determine, where meaningful, whether exact or approximate mixed optimin admits fixed-parameter tractability with respect to $n$, $\max_i m_i$, graph width, payoff rank, or a specified structural parameter. The finite two-person zero-sum case, where mixed optimin coincides with Nash equilibrium and can be computed by linear programming, is a settled boundary case.

\item \textbf{Sensitivity to permitted deviations.}
Repeat the exact-search, approximation, threshold, and verification questions after replacing each $B_j(p)$, consistently throughout the performance function, by either
\[
\widehat B_j^{\mathrm{weak}}(p)
=\{q_j\in\Delta_j:u_j(q_j,p_{-j})\ge u_j(p)\},
\]
\[
\widehat B_j^{\mathrm{best}}(p)
=\{p_j\}\cup\argmax_{q_j\in\Delta_j}u_j(q_j,p_{-j}).
\]
Both induce \emph{modified} performance functions and modified Pareto-undominated sets; they are not definitions of the benchmark strict-deviation optimin. Identify which complexity changes are caused by strict-profitability tests, boundary degeneracies, and simultaneous products of individually defined deviation sets.
\end{enumerate}

\subsection*{Interpretation and scope}

The questions concern \emph{finding one} optimin point, not necessarily enumerating the entire Pareto frontier. Exact mixed-strategy search, a rational $\varepsilon$-approximation to an exact point, additive performance efficiency, weighted-value computation, and global membership testing are different computational tasks. Complexity for a succinctly encoded pure game cannot be inferred from polynomial time for the explicitly tabulated version. Two-player zero-sum games must not be used to infer tractability of arbitrary two-player general-sum games; likewise, the fact that every Nash equilibrium of an $n$-person constant-sum game is optimin does not by itself transfer any hardness theorem in the opposite direction. An algorithm for approximate mixed optimin without a worst-case running-time or approximation theorem does not resolve these classifications.

\subsection*{Source audit and status}

Ismail (2025, Definitions 1--3, Theorem 1, Remark 1, Remark 5, and Proposition 2/Corollary 1) supplies the optimin definition, existence, the pure-restricted benchmark, and the two-person zero-sum identification. It expressly does \emph{not} infer the complexity of mixed optimin from the hardness of computing Nash equilibria. The author's \emph{Optimin-Solver} repository advertises algorithms for pure optimin and approximate mixed optimin in two-player games, but its existence does not establish any worst-case complexity classification. Chen--Deng--Teng (2009), Daskalakis--Goldberg--Papadimitriou (2009), and Etessami--Yannakakis (2010) classify distinct \emph{Nash} computation tasks; their $\mathsf{PPAD}$ and $\mathsf{FIXP}$ results are comparative benchmarks, not theorems about optimin. \`Alvarez--Gabarro--Serna (2011) motivates careful separation of explicit and succinct game representations; it does not classify optimin computation.

The eight subquestions above are newly specified research targets, not eight individually verified open conjectures from the cited papers. On the reviewed sources, a tight worst-case classification for the general mixed-strategy optimin search problem is not established. The explicit-table pure case and the two-person zero-sum case are excluded from the unresolved core. No proof, reduction, or algorithm is asserted in this entry.

\subsection*{References}

\begin{enumerate}
\item[S1] Mehmet S. Ismail (2025). \emph{Super-Nash Performance}. International Economic Review 66(4), 1487--1503. \url{https://doi.org/10.1111/iere.12764}. \textit{Locator:} Section 2, Definitions 1--3 (pp.~1488--1489); Section 3, Theorem 1 and Remark 1 (pp.~1490--1491), Remark 5 and Proposition 2/Corollary 1 (pp.~1492--1493). \textit{Establishes:} the benchmark solution concept and existence/zero-sum properties. \textit{Does not establish:} polynomial-time, hardness, or completeness results for general mixed optimin computation.

\item[S2] Mehmet S. Ismail. \emph{Optimin-Solver} (public software repository). \url{https://github.com/drmehmetismail/Optimin-Solver}. \textit{Locator:} README; \texttt{optimin\_pure.py} and \texttt{optimin\_mixed\_approximate.py}. \textit{Establishes:} availability of implementation examples for pure and approximate mixed optimin in two-person games. \textit{Does not establish:} a worst-case complexity classification.

\item[S3] Xi Chen, Xiaotie Deng and Shang-Hua Teng (2009). \emph{Settling the Complexity of Computing Two-Player Nash Equilibria}. Journal of the ACM 56(3), Article 14. \url{https://doi.org/10.1145/1516512.1516516}. \textit{Locator:} central $\mathsf{PPAD}$-completeness theorem. \textit{Establishes:} complexity of two-player Nash equilibrium. \textit{Does not establish:} the same complexity for two-player optimin.

\item[S4] Constantinos Daskalakis, Paul W. Goldberg and Christos H. Papadimitriou (2009). \emph{The Complexity of Computing a Nash Equilibrium}. SIAM Journal on Computing 39(1), 195--259. \url{https://doi.org/10.1137/070699652}. \textit{Locator:} abstract and principal $\mathsf{PPAD}$-completeness statements. \textit{Establishes:} Nash-computation complexity. \textit{Does not establish:} optimin-computation complexity.

\item[S5] Kousha Etessami and Mihalis Yannakakis (2010). \emph{On the Complexity of Nash Equilibria and Other Fixed Points}. SIAM Journal on Computing 39(6), 2531--2597. \url{https://doi.org/10.1137/080720826}. \textit{Locator:} abstract and the definition of $\mathsf{FIXP}$; exact versus near-equilibrium distinctions. \textit{Establishes:} real-algebraic search benchmarks for Nash equilibria. \textit{Does not establish:} $\mathsf{FIXP}$ membership or hardness for optimin.

\item[S6] Carme \`Alvarez, Joaquim Gabarro and Maria Serna (2011). \emph{Equilibria Problems on Games: Complexity versus Succinctness}. Journal of Computer and System Sciences 77(6), 1172--1197. \url{https://doi.org/10.1016/j.jcss.2011.01.001}. \textit{Locator:} abstract; comparisons of explicit and succinct representations of games. \textit{Establishes:} representation-sensitive equilibrium complexity in other solution concepts. \textit{Does not establish:} the corresponding complexity of optimin.
\end{enumerate}

\subsection*{Common conventions: Source mathematical conventions}

\textbf{Finite games.} For a finite set $A$, $\Delta(A)$ is its probability
simplex. Players choose independent mixed strategies in Nash equilibrium;
correlation is allowed only when explicitly part of the solution concept.
In a game with payoff functions $u_i$ and strategy profile $x$, an additive
$\varepsilon$ Nash equilibrium satisfies
\[
 u_i(x)\geq u_i(x_i',x_{-i})-\varepsilon
 \quad\text{for every player $i$ and every admissible deviation $x_i'$}.
\]
For numerical approximation, payoffs are normalized as locally specified.
An $\varepsilon$ payoff residual is distinct from distance at most
$\varepsilon$ to the set of exact equilibria. Point convergence, convergence
of empirical play, and convergence in distance to an equilibrium set are
also distinct.

\textbf{Computational representations.} Unless stated otherwise, a complexity
question takes rational data with binary encoding; polynomial time is measured
in total bit length. A fixed-accuracy polynomial algorithm is not necessarily
polynomial in $1/\varepsilon$, and neither is necessarily polynomial in
$\log(1/\varepsilon)$. Succinct games and value, demand, or sampling oracles
must declare their interfaces. Exact real solutions require an explicit
representation; an unspecified list of arbitrary real numbers is not a
finite computational output. A claimed algorithmic barrier is meaningful
only for its specified model and reductions.

\textbf{Dynamic games.} Histories, information, observation, and allowable
randomization are part of the model. A stationary strategy depends only on
the current state; a positional strategy in a deterministic graph game
chooses one action at each controlled vertex. Nash equilibrium at an initial
state and equilibrium after every history are different requirements.
An expectation of a pathwise $\liminf$ payoff, a $\liminf$ of expected averages,
a limiting expected average, and a uniform finite-horizon equilibrium are
different criteria. None is silently substituted for another.

\textbf{Allocation and mechanisms.} A complete allocation assigns every item
exactly once unless otherwise stated. Zero-valued goods, negative values,
capacity constraints, feasibility, lotteries, and copies of goods affect
fairness definitions and are specified locally. Dominant-strategy incentive
compatibility, Bayesian incentive compatibility, universal truthfulness,
and truthfulness in expectation impose different inequalities. Whenever
an objective ratio has zero denominator, its convention is stated or those
instances are excluded.

\textbf{Infinite-dimensional models.} Expectations, integrals, stochastic
equations, and optimization objectives require their local regularity,
integrability, filtrations, and admissible domains. A backward--forward
mean-field system is a boundary-value problem; it is not an initial-value
semiflow without an additional construction. Long-time, large-population,
and vanishing-noise limits require an order of limits and a convergence
topology. If a minimum need not be attained, the objective is its infimum
or supremum and the approximation requirement is stated separately.
% END ECONOMICS STATEMENT
\end{document}
